% GENERATED FILE. Edit canonical lessons, then run npm run generate:books.
% canonical-source-hash: fnv1a64:186dc221a171bd36
% canonical-lessons: TA-C470-vicuvacam, TA-C470-turokam, TA-C470-tiyakam, TA-C470-mannippu, TA-C470-paramukam

\chapter{Loyalty and Betrayal}
\label{ch:ta-loyalty-and-betrayal}
\hlchaptermodality{470}

\begin{chapteropening}
\textbf{By the end of this chapter:} \emph{I can talk about loyalty, betrayal, sacrifice and forgiveness.}

Complete the last lesson of chapter 470: I can talk about loyalty, betrayal, sacrifice and forgiveness.
\end{chapteropening}

\section[\texorpdfstring{\ta{விசுவாசம்} (vicuvācam)}{விசுவாசம் (vicuvācam)}]{\ta{விசுவாசம்} (vicuvācam) — loyalty; faithfulness}

\label{lesson:TA-C470-vicuvacam}



\begin{quote}
Before the new one: say the Tamil for an enemy, then the Tamil for to take revenge.
\end{quote}

\subsection*{You'll want to know: \ta{விசுவாசம்}}
\textbf{\ta{விசுவாசம்}} — \emph{vicuvācam} — \textquotedblleft{}loyalty; faithfulness\textquotedblright{}.

\begin{quote}
\ta{அவர்} \ta{இருபது} \ta{வருடமாக} \ta{அந்த} \ta{நிறுவனத்துக்கு} \ta{விசுவாசமாக} \ta{இருந்தார்}.
\end{quote}

\begin{quote}
\emph{avar irupatu varuṭamāka anta niṟuvaṉattukku vicuvācamāka iruntār.}
\end{quote}

\begin{quote}
He was loyal to that company for twenty years.
\end{quote}

\ta{விசுவாசமாக} \ta{இரு}, to be loyal; the one you are loyal to takes -\ta{க்கு}.

\begin{cousinweb}
\textbf{\ta{விசுவாசம்}} (\emph{vicuvācam}) is borrowed from Sanskrit \emph{viśvāsa}, \textquotedblleft{}trust, confidence\textquotedblright{} (CDIAL 11966).
\end{cousinweb}

\begin{grammarlens}[title={What you've built}]
One more word for describing what you lived and felt.
\end{grammarlens}

\subsection*{Your turn}
\begin{itemize}
\raggedright
  \item \emph{Say it:} \emph{vicuvācam}
  \item \emph{Say it:} \emph{vicuvācam}, once more
  \item \emph{Say it:} \emph{avar irupatu varuṭamāka anta niṟuvaṉattukku vicuvācamāka iruntār.} — He was loyal to that company for twenty years.
  \item \emph{Recall:} say the Tamil for an enemy, then the Tamil for to take revenge, then say \emph{vicuvācam} again
\end{itemize}

\begin{tcolorbox}[breakable,colback=teal!4,colframe=teal!35!black,title={Before you move on}]
What does \ta{விசுவாசம்} mean? (\textquotedblleft{}Loyalty; faithfulness\textquotedblright{}.) Say it once more.
\end{tcolorbox}
\section[\texorpdfstring{\ta{துரோகம்} (turōkam)}{துரோகம் (turōkam)}]{\ta{துரோகம்} (turōkam) — betrayal}

\label{lesson:TA-C470-turokam}



\begin{quote}
Before the new one: say the Tamil for to take revenge, then the Tamil for loyalty.
\end{quote}

\subsection*{You'll want to know: \ta{துரோகம்}}
\textbf{\ta{துரோகம்}} — \emph{turōkam} — \textquotedblleft{}betrayal\textquotedblright{}.

\begin{quote}
\ta{அவர்} \ta{நம்பிய} \ta{நண்பனே} \ta{அவருக்குத்} \ta{துரோகம்} \ta{செய்தான்}.
\end{quote}

\begin{quote}
\emph{avar nampiya naṇpaṉē avarukkut turōkam ceytāṉ.}
\end{quote}

\begin{quote}
The very friend he trusted betrayed him.
\end{quote}

\ta{துரோகம்} \ta{செய்}, to betray; the one betrayed takes -\ta{க்கு}.

\begin{cousinweb}
\textbf{\ta{துரோகம்}} (\emph{turōkam}) is borrowed from Sanskrit \emph{drōha}, \textquotedblleft{}injury, treachery\textquotedblright{} (Monier-Williams); Tamil opens the \emph{dr} with a vowel.
\end{cousinweb}

\begin{grammarlens}[title={What you've built}]
One more word for describing what you lived and felt.
\end{grammarlens}

\subsection*{Your turn}
\begin{itemize}
\raggedright
  \item \emph{Say it:} \emph{turōkam}
  \item \emph{Say it:} \emph{turōkam}, once more
  \item \emph{Say it:} \emph{avar nampiya naṇpaṉē avarukkut turōkam ceytāṉ.} — The very friend he trusted betrayed him.
  \item \emph{Recall:} say the Tamil for to take revenge, then the Tamil for loyalty, then say \emph{turōkam} again
\end{itemize}

\begin{tcolorbox}[breakable,colback=teal!4,colframe=teal!35!black,title={Before you move on}]
What does \ta{துரோகம்} mean? (\textquotedblleft{}Betrayal\textquotedblright{}.) Say it once more.
\end{tcolorbox}
\section[\texorpdfstring{\ta{தியாகம்} (tiyākam)}{தியாகம் (tiyākam)}]{\ta{தியாகம்} (tiyākam) — sacrifice}

\label{lesson:TA-C470-tiyakam}



\begin{quote}
Before the new one: say the Tamil for loyalty, then the Tamil for betrayal.
\end{quote}

\subsection*{You'll want to know: \ta{தியாகம்}}
\textbf{\ta{தியாகம்}} — \emph{tiyākam} — \textquotedblleft{}sacrifice\textquotedblright{}.

\begin{quote}
\ta{குழந்தைகளுக்காக} \ta{அம்மா} \ta{பல} \ta{தியாகங்கள்} \ta{செய்தாள்}.
\end{quote}

\begin{quote}
\emph{kuḻantaikaḷukkāka ammā pala tiyākaṅkaḷ ceytāḷ.}
\end{quote}

\begin{quote}
The mother made many sacrifices for her children.
\end{quote}

Plural: \ta{தியாகங்கள்}, \emph{tiyākaṅkaḷ}.

A -\ta{அம்} noun takes -\ta{ங்கள்} in the plural.

\begin{cousinweb}
\textbf{\ta{தியாகம்}} (\emph{tiyākam}) is borrowed from Sanskrit \emph{tyāga}, \textquotedblleft{}giving up; liberality\textquotedblright{} (CDIAL 5985).
\end{cousinweb}

\begin{grammarlens}[title={What you've built}]
One more word for describing what you lived and felt.
\end{grammarlens}

\subsection*{Your turn}
\begin{itemize}
\raggedright
  \item \emph{Say it:} \emph{tiyākam}
  \item \emph{Say it:} \emph{tiyākam}, once more
  \item \emph{Say it:} \emph{kuḻantaikaḷukkāka ammā pala tiyākaṅkaḷ ceytāḷ.} — The mother made many sacrifices for her children.
  \item \emph{Recall:} say the Tamil for loyalty, then the Tamil for betrayal, then say \emph{tiyākam} again
\end{itemize}

\begin{tcolorbox}[breakable,colback=teal!4,colframe=teal!35!black,title={Before you move on}]
What does \ta{தியாகம்} mean? (\textquotedblleft{}Sacrifice\textquotedblright{}.) Say it once more.
\end{tcolorbox}
\section[\texorpdfstring{\ta{மன்னிப்பு} (maṉṉippu)}{மன்னிப்பு (maṉṉippu)}]{\ta{மன்னிப்பு} (maṉṉippu) — forgiveness; an apology}

\label{lesson:TA-C470-mannippu}



\begin{quote}
Before the new one: say the Tamil for betrayal, then the Tamil for sacrifice.
\end{quote}

\subsection*{You'll want to know: \ta{மன்னிப்பு}}
\textbf{\ta{மன்னிப்பு}} — \emph{maṉṉippu} — \textquotedblleft{}forgiveness; an apology\textquotedblright{}.

\begin{quote}
\ta{தவறுக்காக} \ta{அவன்} \ta{என்னிடம்} \ta{மன்னிப்புக்} \ta{கேட்டான்}.
\end{quote}

\begin{quote}
\emph{tavaṟukkāka avaṉ eṉṉiṭam maṉṉippuk kēṭṭāṉ.}
\end{quote}

\begin{quote}
He asked my forgiveness for the mistake.
\end{quote}

\ta{மன்னிப்புக்} \ta{கேள்}, to apologise. The polite \ta{மன்னிக்கவும்} \textquotedblleft{}sorry\textquotedblright{} is the same verb.

\begin{cousinweb}
\textbf{\ta{மன்னிப்பு}} (\emph{maṉṉippu}) is a Tamil noun from the verb \ta{மன்னி} \textquotedblleft{}to forgive\textquotedblright{}; the origin of the verb is not settled in our sources.
\end{cousinweb}

\begin{grammarlens}[title={What you've built}]
One more word for describing what you lived and felt.
\end{grammarlens}

\subsection*{Your turn}
\begin{itemize}
\raggedright
  \item \emph{Say it:} \emph{maṉṉippu}
  \item \emph{Say it:} \emph{maṉṉippu}, once more
  \item \emph{Say it:} \emph{tavaṟukkāka avaṉ eṉṉiṭam maṉṉippuk kēṭṭāṉ.} — He asked my forgiveness for the mistake.
  \item \emph{Recall:} say the Tamil for betrayal, then the Tamil for sacrifice, then say \emph{maṉṉippu} again
\end{itemize}

\begin{tcolorbox}[breakable,colback=teal!4,colframe=teal!35!black,title={Before you move on}]
What does \ta{மன்னிப்பு} mean? (\textquotedblleft{}Forgiveness; an apology\textquotedblright{}.) Say it once more.
\end{tcolorbox}
\section[\texorpdfstring{\ta{பாராமுகம்} (pārāmukam)}{பாராமுகம் (pārāmukam)}]{\ta{பாராமுகம்} (pārāmukam) — the cold shoulder}

\label{lesson:TA-C470-paramukam}



\begin{quote}
Before the new one: say the Tamil for sacrifice, then the Tamil for forgiveness.
\end{quote}

\subsection*{You'll want to know: \ta{பாராமுகம்}}
\textbf{\ta{பாராமுகம்}} — \emph{pārāmukam} — \textquotedblleft{}the cold shoulder\textquotedblright{}.

\begin{quote}
\ta{சண்டைக்குப்} \ta{பிறகு} \ta{அவள்} \ta{என்னிடம்} \ta{பாராமுகமாக} \ta{இருந்தாள்}.
\end{quote}

\begin{quote}
\emph{caṇṭaikkup piṟaku avaḷ eṉṉiṭam pārāmukamāka iruntāḷ.}
\end{quote}

\begin{quote}
After the quarrel she gave me the cold shoulder.
\end{quote}

\ta{பாரா} \textquotedblleft{}not looking\textquotedblright{} + \ta{முகம்} \textquotedblleft{}face\textquotedblright{}.

\begin{cousinweb}
\textbf{\ta{பாராமுகம்}} (\emph{pārāmukam}) is a native compound: the negative of \ta{பார்} \emph{pār}, \textquotedblleft{}to see\textquotedblright{} (DEDR 4091), with \ta{முகம்} \emph{mukam}, \textquotedblleft{}face\textquotedblright{} (DEDR 4889).
\end{cousinweb}

\begin{grammarlens}[title={What you've built}]
One more word for describing what you lived and felt.
\end{grammarlens}

\subsection*{Your turn}
\begin{itemize}
\raggedright
  \item \emph{Say it:} \emph{pārāmukam}
  \item \emph{Say it:} \emph{pārāmukam}, once more
  \item \emph{Say it:} \emph{caṇṭaikkup piṟaku avaḷ eṉṉiṭam pārāmukamāka iruntāḷ.} — After the quarrel she gave me the cold shoulder.
  \item \emph{Recall:} say the Tamil for sacrifice, then the Tamil for forgiveness, then say \emph{pārāmukam} again
\end{itemize}

\begin{tcolorbox}[breakable,colback=teal!4,colframe=teal!35!black,title={Before you move on}]
What does \ta{பாராமுகம்} mean? (\textquotedblleft{}The cold shoulder\textquotedblright{}.) Say it once more.
\end{tcolorbox}
